\documentclass[11pt]{article}
\usepackage[margin=1in]{geometry}
\usepackage{amsmath,amssymb,amsthm,mathtools}
\usepackage[T1]{fontenc}
\usepackage[utf8]{inputenc}
\usepackage{enumitem}
\usepackage{booktabs,longtable,array}
\usepackage{hyperref}
\usepackage{xcolor}
\usepackage{microtype}
\hypersetup{colorlinks=true, linkcolor=blue!50!black, citecolor=blue!50!black, urlcolor=blue!50!black}

\newtheorem{theorem}{Theorem}[section]
\newtheorem{proposition}[theorem]{Proposition}
\newtheorem{lemma}[theorem]{Lemma}
\newtheorem{corollary}[theorem]{Corollary}
\newtheorem{claim}[theorem]{Claim}
\theoremstyle{definition}
\newtheorem{definition}[theorem]{Definition}
\newtheorem{assumption}[theorem]{Assumption}
\theoremstyle{remark}
\newtheorem{remark}[theorem]{Remark}

\newcommand{\KL}{\mathrm{KL}}
\newcommand{\Ent}{\mathrm{Ent}}
\newcommand{\Prob}{\mathsf{P}}
\newcommand{\E}{\mathbb{E}}
\newcommand{\R}{\mathbb{R}}
\newcommand{\N}{\mathbb{N}}
\newcommand{\Pcal}{\mathcal{P}}
\newcommand{\Mcal}{\mathcal{M}}
\newcommand{\Qcal}{\mathcal{Q}}
\newcommand{\Zcal}{\mathcal{Z}}
\newcommand{\Fcal}{\mathcal{F}}
\newcommand{\Gcal}{\mathcal{G}}
\newcommand{\W}{\mathcal{W}}
\newcommand{\one}{\mathbf{1}}
\newcommand{\dd}{\,d}
\newcommand{\supp}{\operatorname{supp}}
\newcommand{\argmin}{\operatorname*{arg\,min}}
\newcommand{\argmax}{\operatorname*{arg\,max}}
\newcommand{\esssup}{\operatorname*{ess\,sup}}
\newcommand{\ri}{\operatorname{ri}}
\newcommand{\dom}{\operatorname{dom}}
\newcommand{\Law}{\operatorname{Law}}
\newcommand{\Var}{\operatorname{Var}}
\newcommand{\Cov}{\operatorname{Cov}}
\newcommand{\diag}{\operatorname{diag}}
\newcommand{\Id}{\mathrm{Id}}
\newcommand{\Lip}{\operatorname{Lip}}
\newcommand{\osc}{\operatorname{osc}}
\newcommand{\diam}{\operatorname{diam}}
\newcommand{\proj}{\operatorname{proj}}
\newcommand{\e}{\mathrm{e}}

\newenvironment{argument}{\paragraph{Argument.}\begin{enumerate}[leftmargin=2em]}{\end{enumerate}}
\newenvironment{proofoutline}{\paragraph{Proof architecture.}\begin{enumerate}[leftmargin=2em]}{\end{enumerate}}


% Source-faithfulness apparatus used by the paper reconstructions.
\newcommand{\sourceanchor}[1]{\par\smallskip\noindent\textbf{Source anchor.} #1\par\smallskip}
\newcommand{\sourcecomment}[1]{\begin{remark}#1\end{remark}}
\newenvironment{sourceguide}
  {\begin{description}[style=nextline,leftmargin=3.2cm,labelwidth=3cm]}
  {\end{description}}

% Backward-compatible environments used by some of the older reconstructions.
\newenvironment{distilled}[1][]{\begin{theorem}[#1]}{\end{theorem}}
\newenvironment{distillation}{\begin{enumerate}[leftmargin=2em]}{\end{enumerate}}

\newenvironment{commentarynotes}{\begin{remark}\begin{enumerate}[leftmargin=2em]}{\end{enumerate}\end{remark}}

\title{Fortet's Fixed-Point Method for the Schr\"odinger System}
\author{}
\date{L\'eonard (2019)}
\begin{document}
\maketitle

\section*{Source map and scope}
\begin{description}[style=nextline,leftmargin=3.2cm,labelwidth=3cm]
\item[Primary source] \texttt{\detokenize{Leonard 2019-revisit-fortet-1904.13211v1.pdf}}
\item[Coverage] The fixed-point proof of Fortet's theorem: reduction, monotone scheme, uniqueness, positivity dichotomy, existence theorem, and the variants in Sections 6--7.
\item[Source anchors]
\begin{itemize}[leftmargin=2em]
  \item Theorems 1.1 and 1.9
  \item Proposition 4.4
  \item Lemmas 4.1--4.13
  \item Proposition 5.1
  \item Theorem 5.6
  \item Lemmas 5.7--5.12
  \item Proposition 6.1
  \item Proposition 7.2
\end{itemize}
\end{description}

The exposition below preserves the source proof order and proof technique.  Notation is normalized
and intermediate calculations are made explicit.  When the source contains a gap, ambiguity, or
apparently incorrect step, the reconstruction records the source step before adding a separate
commentary remark.

\section{The Schr\"odinger system}
Let $(X,\mathcal X)$ and $(Y,\mathcal Y)$ be measurable spaces, let
$\mu\in\mathcal P(X)$ and $\nu\in\mathcal P(Y)$, and let
$p:X\times Y\to(0,\infty)$ be measurable.  The Schr\"odinger system asks for
$\sigma$-finite positive measures $a$ and $b$ such that
\begin{align}
 a(dx)\int_Yp(x,y)b(dy)&=\mu(dx),\label{eq:L19-sch1}\\
 b(dy)\int_Xp(x,y)a(dx)&=\nu(dy).\label{eq:L19-sch2}
\end{align}
Only the product measure
\[
 \pi(dx,dy)=p(x,y)a(dx)b(dy)
\]
is intrinsic: $(a,b)$ and $(c^{-1}a,cb)$ determine the same solution for every $c>0$.

Write $u=d\mu/da$.  Then
\[
 a=u^{-1}\mu,
 \qquad
 b(dy)=\Psi[u](y)^{-1}\nu(dy),
\]
where
\begin{equation}\label{eq:L19-Psi}
 \Psi[u](y):=\int_Xp(x,y)u(x)^{-1}\,\mu(dx).
\end{equation}
Substitution into the first equation yields Fortet's fixed-point equation
\begin{equation}\label{eq:L19-fixed}
 u=\Phi[u],
 \qquad
 \Phi[u](x):=\int_Yp(x,y)\Psi[u](y)^{-1}\,\nu(dy).
\end{equation}
The conventions $0^{-1}=\infty$ and $\infty^{-1}=0$ are used, with products interpreted
through monotone truncation.

\section{Order and homogeneity}
The two reciprocal integral maps reverse order separately:
\[
 u\le v\quad\Longrightarrow\quad
 \Psi[u]\ge\Psi[v]
 \quad\Longrightarrow\quad
 \Phi[u]\le\Phi[v].
\]
Thus $\Phi$ is order preserving.  It is positively homogeneous:
\begin{equation}\label{eq:L19-homogeneous}
 \Psi[cu]=c^{-1}\Psi[u],
 \qquad
 \Phi[cu]=c\Phi[u]
 \quad(c>0).
\end{equation}
This is the projective symmetry of the Schr\"odinger system.

Whenever the expressions are finite on the relevant sets, Tonelli's theorem gives the
mass identity
\begin{equation}\label{eq:L19-mass-id}
 \int_{\{u>0\}}\frac{\Phi[u]}{u}\,d\mu
 =\nu\bigl(\Psi[u]<\infty\bigr).
\end{equation}
Indeed,
\begin{align*}
 \int\frac{\Phi[u](x)}{u(x)}\,\mu(dx)
 &=\int_Y\Psi[u](y)^{-1}
    \left(\int_Xp(x,y)u(x)^{-1}\,\mu(dx)\right)\nu(dy)\\
 &=\int_Y\mathbf1_{\{\Psi[u]<\infty\}}\,\nu(dy).
\end{align*}
When $u$ is strictly positive and comparable to an admissible envelope, both sides equal
one.

\section{The monotone Fortet iteration}
Choose an envelope $U:X\to(0,\infty)$ such that $\Psi[U]$ and $\Phi[U]$ are finite.
Fortet's regularized iteration is
\begin{equation}\label{eq:L19-iteration}
 u_1=U,
 \qquad
 u_{n+1}=\frac{U}{n+1}\vee\bigl(\Phi[u_n]\wedge U\bigr).
\end{equation}
The floor $U/(n+1)$ keeps every iterate positive; the ceiling $U$ gives a common bound.
Monotonicity of $\Phi$ yields
\begin{equation}\label{eq:L19-monotone}
 0<u_{n+1}\le u_n\le U,
 \qquad
 \Phi[u_{n+1}]\le\Phi[u_n].
\end{equation}
Hence $u_n\downarrow u_*$ pointwise for some $0\le u_*\le U$.
Monotone convergence through the reciprocal maps gives
\begin{equation}\label{eq:L19-truncated-fixed}
 u_*=\Phi[u_*]\wedge U
 \quad\mu\text{-almost everywhere}.
\end{equation}
The remaining problem is to prove that $\Phi[u_*]$ is positive.  Once positivity is known,
the mass identity removes the truncation.

\section{From positivity to an exact fixed point}
Assume $\Phi[u_*]>0$ $\mu$-almost everywhere.  On the set where
$\Phi[u_*]\le U$, equation \eqref{eq:L19-truncated-fixed} gives
$\Phi[u_*]/u_*=1$; on the complementary set the ratio is strictly larger than one.
Applying \eqref{eq:L19-mass-id} and using that both probability masses equal one forces
the complementary set to have zero measure.  Therefore
\begin{equation}\label{eq:L19-ae-fixed}
 0<u_*=\Phi[u_*]\le U
 \quad\mu\text{-almost everywhere}.
\end{equation}

To upgrade an almost-everywhere identity to an everywhere identity, note that changing
$u$ on a $\mu$-null set leaves $\Psi[u]$ and therefore $\Phi[u]$ unchanged everywhere.
Replacing $u_*$ by the everywhere-defined function $\Phi[u_*]$ gives
\begin{equation}\label{eq:L19-everywhere-fixed}
 0<u_*(x)=\Phi[u_*](x)\le U(x)
 \qquad(x\in X).
\end{equation}

\section{The positivity argument under integral hypotheses}
The first principal theorem assumes
\begin{enumerate}
\item $p(x,y)>0$ for all $(x,y)$;
\item $p$ is bounded above;
\item at least one of
\begin{align}
 \int_Y\left(\int_Xp(x,y)\,\mu(dx)\right)^{-1}\nu(dy)&<\infty,
 \label{eq:L19-int1}\\
 \int_X\left(\int_Yp(x,y)\,\nu(dy)\right)^{-1}\mu(dx)&<\infty
 \label{eq:L19-int2}
\end{align}
 holds.
\end{enumerate}
By symmetry it suffices to use \eqref{eq:L19-int1} and take $U=1$.

Strict positivity of $p$ gives a dichotomy:
\begin{equation}\label{eq:L19-dichotomy}
 \text{for }0\le u\le U,
 \qquad
 \Phi[u]\equiv0
 \quad\text{or}\quad
 \Phi[u](x)>0\text{ for every }x.
\end{equation}
Indeed, if $\Phi[u](x_0)=0$, then positivity of $p(x_0,\cdot)$ forces
$\Psi[u]^{-1}=0$ $\nu$-almost everywhere, hence $\Phi[u]\equiv0$.

Suppose for contradiction that $\Phi[u_*]\equiv0$.  Put
\[
 \gamma_u(dy)=\Psi[u](y)^{-1}\nu(dy).
\]
Since $u_n\downarrow u_*$, the measures satisfy
$0\le\gamma_{u_{n+1}}\le\gamma_{u_n}\le\gamma_{u_1}$, and
\[
 \Phi[u_n](x)=\int p(x,y)\,\gamma_{u_n}(dy)\downarrow0.
\]
Fix $x_0$.  On $Y_k=\{y:p(x_0,y)\ge1/k\}$,
\[
 \gamma_{u_n}(Y_k)
 \le k\int p(x_0,y)\,\gamma_{u_n}(dy)
 =k\Phi[u_n](x_0)\longrightarrow0.
\]
Because $p(x_0,y)>0$, $Y_k\uparrow Y$, and the integrability assumption makes
$\gamma_{u_1}$ finite.  Thus the mass of $Y\setminus Y_k$ is uniformly small for large
$k$.  Boundedness of $p$ then gives
\[
 \sup_x\Phi[u_n](x)
 \le \|p\|_\infty
   \bigl(\gamma_{u_n}(Y_k)+\gamma_{u_1}(Y\setminus Y_k)\bigr)
 \longrightarrow0.
\]
Hence for some $n_0$, $\Phi[u_{n_0}]\le U$.  The mass identity and the definition of the
next iterate then force
\[
 u_{n_0+1}=\Phi[u_{n_0+1}]>0,
\]
and monotonicity implies $u_*=u_{n_0+1}$, contradicting
$\Phi[u_*]\equiv0$.  Therefore $\Phi[u_*]>0$ everywhere, and
\eqref{eq:L19-everywhere-fixed} follows.

\begin{theorem}[Existence under Fortet's integral criterion]
If $p$ is positive and bounded and either \eqref{eq:L19-int1} or
\eqref{eq:L19-int2} holds, then the Schr\"odinger system has a solution.
The associated coupling $p(x,y)a(dx)b(dy)$ is unique, and the pair $(a,b)$ is unique up
to reciprocal scaling.
\end{theorem}

\section{Uniqueness of the fixed ray}
Let $u'$ and $u''$ be two strictly positive fixed points.  By homogeneity, normalize them
at some $x_0$ so that $u'(x_0)=u''(x_0)$, and define
\[
 \alpha=\sup\{c>0:cu'\le u''\},
 \qquad
 \beta=\inf\{c>0:u''\le cu'\}.
\]
Order preservation gives $\alpha u'\le u''\le\beta u'$.  At a touching point, equality of
$\Phi[u']$ and $\Phi[u'']$ together with strict positivity of $p$ forces equality of the
intermediate reciprocal integrals $\Psi[u']$ and $\Psi[u'']$, hence equality of the fixed
points after normalization.  Therefore every positive fixed point lies on the same ray.
Equivalently, the coupling is unique and the scaling measures have only the expected
one-dimensional gauge freedom.

\section{Topological variant}
The integral criterion can be replaced by a compact-uniform-integrability condition.
Assume $X$ is topological and there is a nonempty compact $K\subset X$ such that
$p(\cdot,y)$ is continuous on $K$ for every $y$, and there exist finitely many points
$x_j$ and constants $c_j>0$ with
\begin{equation}\label{eq:L19-domination}
 \sup_{x\in K}p(x,y)\le\sum_{j=1}^Jc_jp(x_j,y)
 \qquad(y\in Y).
\end{equation}
The domination in \eqref{eq:L19-domination} makes the family of kernel integrands
uniformly integrable on $K$.  If $\Phi[u_*]$ vanished at one point, the same argument as
above would give pointwise convergence $\Phi[u_n]\to0$; continuity and uniform
integrability upgrade this to uniform convergence on $K$.  A multiplicative change of
kernel then transports the bound from $K$ to all of $X$, providing an iterate below the
envelope and closing the positivity argument.

\begin{theorem}[Existence under compact kernel control]
If $p$ is positive and bounded and satisfies the compact continuity and finite domination
condition \eqref{eq:L19-domination}, then the Schr\"odinger system has a unique coupling.
\end{theorem}

For a translation kernel $p(x,y)=\theta(|y-x|)$ on $\mathbb R^d$, positivity,
boundedness, continuity, and eventual monotonicity of $\theta$ imply the finite
domination condition on every compact set, yielding existence and uniqueness.

\section{Recovery of the scaling measures}
Given the positive fixed point $u_*$, define
\[
 a(dx)=u_*(x)^{-1}\mu(dx),
 \qquad
 b(dy)=\Psi[u_*](y)^{-1}\nu(dy).
\]
Then
\[
 \int_Yp(x,y)b(dy)=\Phi[u_*](x)=u_*(x),
\]
which proves \eqref{eq:L19-sch1}, and the definition of $b$ proves
\eqref{eq:L19-sch2}.  Thus the monotone scalar iteration constructs the full
Schr\"odinger factorization.

\begin{thebibliography}{99}
\bibitem{ChenGeorgiouPavon2021} Yongxin Chen, Tryphon T. Georgiou, and Michele Pavon. \emph{Stochastic Control Liaisons: Richard Sinkhorn Meets Gaspard Monge on a Schr\"odinger Bridge}. SIAM Review, 63(2):249--313, 2021. doi:10.1137/20M1339982. arXiv:2005.10963.
\bibitem{Leonard2014Survey} Christian L\'eonard. \emph{A Survey of the Schr\"odinger Problem and Some of Its Connections with Optimal Transport}. Discrete and Continuous Dynamical Systems--A, 34(4):1533--1574, 2014. doi:10.3934/dcds.2014.34.1533. arXiv:1308.0215.
\bibitem{Fortet1940} Robert Fortet. \emph{Resolution d'un systeme d'equations de M. Schrodinger}. Journal de mathematiques pures et appliquees, 9e serie, 19(1--4):83--105, 1940.
\bibitem{SinkhornKnopp1967} Richard Sinkhorn and Paul Knopp. \emph{Concerning Nonnegative Matrices and Doubly Stochastic Matrices}. Pacific Journal of Mathematics, 21(2):343--348, 1967.
\bibitem{FranklinLorenz1989} Joel Franklin and Jens Lorenz. \emph{On the Scaling of Multidimensional Matrices}. Linear Algebra and its Applications, 114/115:717--735, 1989. doi:10.1016/0024-3795(89)90490-4.
\bibitem{Carroll2004} Joseph E. Carroll. \emph{Birkhoff's Contraction Coefficient}. Linear Algebra and its Applications, 389:227--234, 2004. doi:10.1016/j.laa.2004.02.039.
\bibitem{EvesonNussbaum1995} Simon P. Eveson and Roger D. Nussbaum. \emph{An Elementary Proof of the Birkhoff--Hopf Theorem}. Mathematical Proceedings of the Cambridge Philosophical Society, 117(1):31--55, 1995.
\bibitem{BlanchetMurthy2019} Jose Blanchet and Karthyek R. A. Murthy. \emph{Quantifying Distributional Model Risk via Optimal Transport}. Mathematics of Operations Research, 44(2):565--600, 2019. doi:10.1287/moor.2018.0936. arXiv:1604.01446.
\bibitem{GaoKleywegt2023} Rui Gao and Anton J. Kleywegt. \emph{Distributionally Robust Stochastic Optimization with Wasserstein Distance}. Mathematics of Operations Research, 48(2):603--655, 2023. doi:10.1287/moor.2022.1275. arXiv:1604.02199.
\bibitem{MohajerinEsfahaniKuhn2018} Peyman Mohajerin Esfahani and Daniel Kuhn. \emph{Data-Driven Distributionally Robust Optimization Using the Wasserstein Metric: Performance Guarantees and Tractable Reformulations}. Mathematical Programming, 171(1--2):115--166, 2018. doi:10.1007/s10107-017-1172-1. arXiv:1505.05116.
\bibitem{WangGaoXie2025} Jie Wang, Rui Gao, and Yao Xie. \emph{Sinkhorn Distributionally Robust Optimization}. Operations Research, 2025. doi:10.1287/opre.2023.0294. arXiv:2109.11926.
\bibitem{Girsanov1960} I. V. Girsanov. \emph{On Transforming a Certain Class of Stochastic Processes by Absolutely Continuous Substitution of Measures}. Theory of Probability and Its Applications, 5(3):285--301, 1960. doi:10.1137/1105027.
\bibitem{CameronMartin1945} R. H. Cameron and W. T. Martin. \emph{Transformations of Wiener Integrals under a General Class of Linear Transformations}. Transactions of the American Mathematical Society, 58(2):184--219, 1945.
\bibitem{Yeh1978} J. Yeh. \emph{Transformation of Conditional Wiener Integrals under Translation and the Cameron--Martin Translation Theorem}. Tohoku Mathematical Journal, 30(4):505--515, 1978.
\bibitem{Kakutani1948} Shizuo Kakutani. \emph{On Equivalence of Infinite Product Measures}. Annals of Mathematics, 49(1):214--224, 1948.
\bibitem{ShiDeBortoliCampbellDoucet2023} Yuyang Shi, Valentin De Bortoli, Andrew Campbell, and Arnaud Doucet. \emph{Diffusion Schr\"odinger Bridge Matching}. Advances in Neural Information Processing Systems, 2023. arXiv:2303.16852.
\bibitem{LiuLipmanNickelKarrerTheodorouChen2024} Guan-Horng Liu, Yaron Lipman, Maximilian Nickel, Brian Karrer, Evangelos A. Theodorou, and Ricky T. Q. Chen. \emph{Generalized Schr\"odinger Bridge Matching}. International Conference on Learning Representations, 2024. arXiv:2310.02233.
\bibitem{DomingoEnrichDrozdzalKarrerChen2025} Carles Domingo-Enrich, Michal Drozdzal, Brian Karrer, and Ricky T. Q. Chen. \emph{Adjoint Matching: Fine-Tuning Flow and Diffusion Generative Models with Memoryless Stochastic Optimal Control}. arXiv preprint arXiv:2409.08861, 2025.

\bibitem{Birkhoff1957} Garrett Birkhoff. \emph{Extensions of Jentzsch's Theorem}. Transactions of the American Mathematical Society, 85(1):219--227, 1957.
\bibitem{Sinkhorn1964} Richard Sinkhorn. \emph{A Relationship Between Arbitrary Positive Matrices and Doubly Stochastic Matrices}. Annals of Mathematical Statistics, 35(2):876--879, 1964.
\bibitem{Sinkhorn1967Prescribed} Richard Sinkhorn. \emph{Diagonal Equivalence to Matrices with Prescribed Row and Column Sums}. American Mathematical Monthly, 74(4):402--405, 1967.
\bibitem{Bauer1965} Friedrich L. Bauer. \emph{An Elementary Proof of the Hopf Inequality for Positive Operators}. Numerische Mathematik, 7:331--337, 1965.
\bibitem{Bushell1973} P. J. Bushell. \emph{Hilbert's Metric and Positive Contraction Mappings in a Banach Space}. Archive for Rational Mechanics and Analysis, 52:330--338, 1973.

\end{thebibliography}

\end{document}
